\begin{abstract}
We classify 150 patients of the ACDC dataset into five diagnostic groups (normal, myocardial infarction, dilated and hypertrophic cardiomyopathy, abnormal right ventricle) from segmentations of their end-diastolic and end-systolic cardiac MRI, height and weight. The 100 training patients are labeled; in the 50 test patients, the left-ventricular (LV) cavity is not segmented and must be reconstructed. With so few patients, the evaluation protocol matters as much as the model. On 21 classical volumetric biomarkers, selecting the hyper-parameters of standard classifiers on a single cross-validation split and reporting the best score gives accuracies of 0.91 to 0.95; nested cross-validation repeated over several splits gives 0.90 to 0.94 for the same models and grids, and an ensemble weight chosen on the evaluation folds varies from 0 to 0.9 with the random seed. Filling the holes of the myocardium recovers the LV cavity with a mean Dice of 0.997; a convex-hull fallback on the few slices where the myocardial ring is open brings it to 0.9998, and training on reconstructed or true cavities makes no measurable difference. What does improve classification is describing the myocardial wall: 11 thickness features raise the accuracy of linear models and of an ensemble of 50 MLPs by 3 to 6 points. The final model, an average of this ensemble and a linear SVM, reaches $0.966 \pm 0.011$ in nested cross-validation; its remaining errors are infarcted hearts that look dilated, and conversely.
\end{abstract}
